\Needspace{14\baselineskip}
\section{Experiment 2: Early Rationale Attribution for Uncertainty Quantification}
\label{sec:era_method}
\markboth{\thesection\quad Experiment 2: ERA Method Design}{}

\subsection{Methodological Motivation}
\label{subsec:era_motivation}

Experiment~1 shows that existing UQ methods, including Perplexity, Semantic Entropy, and UMPIRE, can distinguish correct from incorrect responses to some extent on the full evaluation data, but fail conspicuously on low-uncertainty hallucination (LUH). In the constructed low-uncertainty subsets, hallucinated positives and non-hallucinated negatives are deliberately matched in the score space of the three baseline methods, where fixed-direction AUROC values are near random or mildly reverse ordered. This result indicates that generation probabilities, consistency across repeated samples, and the stability of final hidden representations are insufficient for identifying cases in which a model produces erroneous content consistently.

The underlying reason is that conventional UQ methods primarily answer the following question: \textbf{Does the model consistently believe its current response?}

LUH exposes a different problem. Even when a model assigns high generative confidence to a response, that response may still lack genuine support from the input image and question \citep{shi2024trusting}. In other words, \textbf{predictive stability is not equivalent to evidential sufficiency}.

Motivated by this observation, Experiment~2 no longer searches for additional uncertainty signals in the output probability distribution. Instead, it examines the \emph{source structure of information} involved in forming the final answer. After format-oriented LoRA fine-tuning, a structured response takes the form
\[
\texttt{<vision>}V\texttt{</vision>}
\texttt{<reasoning>}R\texttt{</reasoning>}
\texttt{<answer>}A\texttt{</answer>}.
\]
When generating the final answer $A$, the model may draw on two qualitatively different sources of information. The image $I$ and question text $Q$ originate outside the model and can therefore be regarded as external evidence for the current problem. By contrast, the visual description $V$ and reasoning process $R$ have already been generated and internally transformed by the model; they may contain perceptual or reasoning errors introduced earlier in the response \citep{zhang2024snowball}.

I therefore propose \textbf{ERA (Early Rationale Attribution)}, based on the following hypothesis:

\begin{gapblock}
\textbf{The more confident a model is in a response, the more its answer generation relies on the model's own previously generated visual description and reasoning, and the less it refers to the image and question. This inward dependence is already visible at the early (shallow) stage of answer formation, so preceding errors are self-reinforced and ultimately manifest as low-uncertainty hallucinations.}
\end{gapblock}

ERA therefore does not measure how unstable an answer is. Instead, it measures the \textbf{dependency bias of the answer decision between self-generated content and external evidence}. This quantity can be interpreted as a structural form of uncertainty grounded in evidence provenance.

\subsection{Five Semantic Regions}
\label{subsec:era_regions}

To characterize the information sources of the final answer, ERA partitions the tokens in a complete forward pass into the five semantic regions shown in \Cref{tab:era_regions}.

\begin{table}[h]
\centering
\caption{Five semantic regions used by ERA.}
\label{tab:era_regions}
\datatablefont
\begin{tabularx}{\linewidth}{L{2.5cm}C{1.4cm}Y L{3.1cm}}
\toprule
\tableheadrow \thead{Region} & \thead{Symbol} & \thead{Content} & \thead{Source type} \\
\midrule
\textbf{Image evidence} & $I$ & Visual tokens corresponding to the input image & External evidence \\
\tablestriperow \textbf{Question prompt} & $Q$ & System prompt, user question, and other prompt-text tokens & External evidence \\
\textbf{Visual description} & $V$ & Content within \texttt{<vision>...\allowbreak</vision>} & Self-generated context \\
\tablestriperow \textbf{Reasoning process} & $R$ & Content within \texttt{<reasoning>...\allowbreak</reasoning>} & Self-generated context \\
\tableaccentrow \textbf{Final answer} & $A$ & Content within \texttt{<answer>...\allowbreak</answer>} & Decision under analysis \\
\bottomrule
\end{tabularx}
\end{table}

Here, $I+Q$ denotes the original problem evidence directly available to the model, whereas $V+R$ denotes intermediate information generated by the model after processing that evidence. This partition lets ERA address the following central question directly: \textbf{When producing its answer, does the model attend more strongly to external evidence or to its own preceding rationale?}

\subsection{Attention Attribution from the Answer to Evidence Regions}
\label{subsec:era_attribution}

ERA uses attention distributions to quantify the dependence of the answer on each region \citep{huang2024opera}. Let $\mathbf A^{(l,h)}$ denote the attention matrix of head $h$ in Transformer decoder layer $l$. In an autoregressive model, the hidden state at position $t-1$ predicts the token at position $t$. Therefore, for an answer token $t\in A$, ERA uses the corresponding prediction row $t-1$, rather than the row occupied by token $t$ itself. This convention ensures that the measured attention corresponds to the decision generating the current answer token and matches the prediction-row definition used in the ERA implementation.

For any target region $T\in\{I,Q,V,R\}$, the average attention mass from the answer to $T$ in layer $l$ is defined as
\[
\alpha_l(A\rightarrow T)
=
\frac{1}{H\lvert A\rvert}
\sum_{t\in A}
\sum_{h=1}^{H}
\sum_{j\in T}
A^{(l,h)}_{t-1,j},
\]
where $H$ is the number of attention heads and $\lvert A\rvert$ is the number of tokens in the answer region.

For each answer token, this quantity first sums its attention over all tokens in the target region and then averages over all answer tokens and attention heads. It consequently measures \textbf{the average amount of attention allocated to a particular type of context while the final answer is being generated.}

The procedure yields four attribution quantities:
\[
\alpha_l(A\rightarrow I),\qquad
\alpha_l(A\rightarrow Q),\qquad
\alpha_l(A\rightarrow V),\qquad
\alpha_l(A\rightarrow R).
\]
The first two capture the answer's direct dependence on external input, whereas the latter two capture its dependence on self-generated intermediate reasoning.

\subsection{ERA Uncertainty Score}
\label{subsec:era_score}

Based on the preceding attribution quantities, the ERA score at layer $l$ is defined as
\[
U_{\mathrm{ERA}}^{(l)}
=
\frac{
\alpha_l(A\rightarrow V)+
\alpha_l(A\rightarrow R)
}{
\alpha_l(A\rightarrow I)+
\alpha_l(A\rightarrow Q)+
\alpha_l(A\rightarrow V)+
\alpha_l(A\rightarrow R)+\epsilon
},
\]
where $\epsilon$ is a small constant used to prevent division by zero. Attention from the answer to itself, $A\rightarrow A$, is excluded from both the numerator and denominator. This component primarily reflects the autoregressive linguistic coherence of the answer and grows systematically with answer length. Excluding it restricts the score to the relative dependence on external evidence and self-generated reasoning, reducing sensitivity to response length.

The sample-level score is the average over the first two decoder layers:
\[
U_{\mathrm{ERA}}
=
\frac{1}{2}
\left(
U_{\mathrm{ERA}}^{(0)}+
U_{\mathrm{ERA}}^{(1)}
\right).
\]
Attention in shallow layers more directly reflects the initial routing of information involved in answer formation \citep{belrose2023eliciting}. As depth increases, information from different sources becomes increasingly mixed, weakening the ability of direct attention to distinguish those sources.

A small $U_{\mathrm{ERA}}$ indicates that the final answer depends more directly on external evidence from the image and question. A large value indicates that the answer is built primarily on the model's own preceding visual description and reasoning.

\subsection{Single-Forward-Pass Computation}
\label{subsec:era_computation}

Unlike Semantic Entropy and UMPIRE, ERA does not require additional rounds of stochastic sampling. Given the greedy response already produced in Experiment~1, ERA computes the score with a single teacher-forced forward replay. The overall pipeline is illustrated in \Cref{fig:era_computation_flow} and proceeds as follows:

\begin{figure}[h]
\centering
\includegraphics[width=0.95\linewidth]{ERA计算流程}
\caption{Single-forward-pass computation flow of the Early Rationale Attribution (ERA) uncertainty score. Given the input $(I, Q)$ and the previously generated greedy response $(V, R, A)$, the model executes a single teacher-forced forward replay to extract attention matrices at early decoder layers ($l \in \{0, 1\}$), accumulates attribution mass from answer prediction rows to semantic regions, and computes the relative rationale attribution score.}
\label{fig:era_computation_flow}
\end{figure}

\begin{enumerate}[label=\textbf{Step \arabic*:}, leftmargin=5.2em, labelsep=0.7em]
\item \textbf{Sequence Reconstruction:} For sample $x=(I,Q)$, retrieve the original token sequence of the structured greedy response $Y=(V,R,A)$ and concatenate $[I,Q,V,R,A]$ into a complete teacher-forced input sequence.
\item \textbf{Single Forward Replay:} Feed the assembled sequence through the model in a single teacher-forced forward pass with the original image $I$ kept unchanged, extracting attention weight matrices $\mathbf{A}^{(l,h)}$ from shallow decoder layers ($l \in \{0, 1\}$).
\item \textbf{Attention Attribution Extraction:} For each answer prediction row $t-1$ ($t \in A$), aggregate attention mass across heads $h$ and target tokens $j \in T$ to obtain the regional attribution $\alpha_l(A \rightarrow T)$ for $T \in \{I, Q, V, R\}$.
\item \textbf{Layerwise Score \& Aggregation:} Compute the internal rationale ratio $U_{\mathrm{ERA}}^{(l)}$ for Layers 0 and 1 via \Cref{subsec:era_score}, and average them to obtain the final sample-level uncertainty score $U_{\mathrm{ERA}} = \frac{1}{2}(U_{\mathrm{ERA}}^{(0)} + U_{\mathrm{ERA}}^{(1)})$.
\end{enumerate}

Compared with consistency-based methods that require repeated sampling, ERA adds only one complete forward replay of the existing greedy response and generates no additional candidate answers.
